%% kapa: exam=1 round=19 year=2008 period=1 form=A subject=civil_law subject_ko=민법 seq=34 number=34
%% source: Q-Net 감정평가사 기출문제 (cst00309) artlSeq=1001962 file=extracted/1008/1001962/2008년 제19회 감정평가사 1차시험 기출문제/1교시.hwp
%% conversion: kordoc 4.24.51 -> tools/gen1.py (auto); review=proofread-sonnet; flags=none
\begin{question}{34}
관습법상의 법정지상권에 관한 설명으로 틀린 것은? (다툼이 있는 경우에는 판례에 의함)
\begin{choices}
  \item 관습법상의 법정지상권이 성립하려면 토지와 건물의 소유권이 분리될 당시에 토지와 그 지상건물이 동일인의 소유에 속하였어야 한다.
  \item 적법하게 명의신탁된 토지위에 수탁자가 건물을 신축한 후 명의신탁이 해지되어 토지소유권이 신탁자에게 환원된 경우, 수탁자는 토지에 대하여 관습법상의 법정지상권을 취득하지 못한다.
  \item 관습법상의 법정지상권은 관습법에 의하여 성립하는 것이므로 등기 없이도 취득한다.
  \item 관습법상의 법정지상권을 취득한 자는 이를 취득할 당시의 토지소유자에게는 등기 없이도 그 지상권을 주장할 수 있지만, 토지소유자로부터 소유권을 양수한 제3자에게는 등기를 하여야만 그 지상권을 주장할 수 있다.
  \item 법정지상권을 가진 건물소유자로부터 건물을 양수하면서 지상권까지 양도받기로 한 자에 대하여 토지소유자가 건물철거 및 대지의 인도를 구하는 것은 신의성실의 원칙상 허용될 수 없다.
\end{choices}
\end{question}
\answer{4}
